\chapter{Resultados --- Fred (base de datos, reproducibilidad y rendimiento)}
\label{cap:resultados-fred}

Este cap\'itulo integra la informaci\'on real disponible en
\texttt{docs/adr/ADR-004-postgresql.md},
\texttt{docs/adr/ADR-005-despliegue.md},
\texttt{docs/basedatos/CATALOGO-SP.md} y
\texttt{docs/mediciones/perf/REPORT.md} sobre el trabajo de Fred. Ninguna
cifra de este cap\'itulo se invent\'o: todas provienen de esos documentos o
de \texttt{docs/mediciones/perf/k6-runN-\{frio,caliente\}.json}.

\section{Base de datos y PostgreSQL}

El esquema de PostgreSQL~16~\cite{postgresql-docs} se gestiona con Flyway
(\texttt{V1__schema_inicial.sql}) como \'unica fuente de verdad para la
aplicaci\'on en tiempo de ejecuci\'on. Para permitir que Docker reconstruya
el sistema de forma id\'entica desde una clonaci\'on limpia sin depender de
que Spring Boot arranque primero, ADR-004 introduce un mecanismo doble:

\begin{itemize}[nosep]
  \item \texttt{db/schema.sql}: r\'eplica exacta del esquema de Flyway,
    montada en \texttt{docker-entrypoint-initdb.d} para que Postgres cree
    las tablas en el primer arranque del contenedor.
  \item \texttt{db/roles.sql}: crea el rol \texttt{biopet_app}
    (\texttt{LOGIN}, sin \texttt{SUPERUSER}/\texttt{CREATEDB}/\texttt{CREATEROLE})
    con privilegios m\'inimos: \texttt{SELECT/INSERT/UPDATE/DELETE}
    \'unicamente sobre \texttt{usuarios} y \texttt{mascotas}.
  \item \texttt{db/seed.sql}: siembra el usuario administrador de forma
    idempotente (\texttt{WHERE NOT EXISTS}).
\end{itemize}

Hibernate/JPA se conecta en tiempo de ejecuci\'on con \texttt{biopet_app};
Flyway sigue usando la cuenta propietaria \texttt{biopet_user}, separando
la conexi\'on que aplica migraciones de la que sirve tr\'afico de la
aplicaci\'on. Gracias a \texttt{spring.flyway.baseline-on-migrate: true},
Flyway reconoce el esquema ya creado por \texttt{db/schema.sql} y establece
un baseline en la versi\'on~1, sin reintentar la migraci\'on original.
Verificado con \texttt{psql}: \texttt{biopet_app} no tiene atributos de
\texttt{Superuser}, \texttt{Create role} ni \texttt{Create DB}; verificado
con \texttt{flyway_schema_history}: una \'unica fila
\texttt{version~1, type=BASELINE, success=true}.

% EVIDENCIA-FRED-02
% Ruta esperada: figuras/fred/02-postgresql-procedimiento.png
\evidencia{figuras/fred/02-postgresql-procedimiento.png}%
{Ejecuci\'on real de \texttt{fn\_resumen\_mascotas\_por\_especie} conectado como \texttt{biopet\_app}, sin GRANT adicional.}%
{fig:fred-postgresql-procedimiento}

\section{Funci\'on de resumen y procedimientos almacenados}

\texttt{docs/basedatos/CATALOGO-SP.md} documenta dos objetos PL/pgSQL:

\begin{itemize}[nosep]
  \item \textbf{\texttt{fn_resumen_mascotas_por_especie}}: funci\'on (no
    procedimiento, para permitir invocaci\'on directa v\'ia
    \texttt{@Query(nativeQuery=true)}), par\'ametro opcional
    \texttt{p_duenio_id BIGINT}, devuelve \texttt{especie}/\texttt{total}
    de mascotas activas. Expuesta por
    \texttt{GET /api/mascotas/resumen-especies}. Verificada con
    Testcontainers (\texttt{ResumenEspeciesIntegrationTest}) y de forma
    manual: \texttt{Perro~|~1} (29/07/2026).
  \item \textbf{\texttt{set_actualizado_en}}: funci\'on \emph{trigger}
    (\texttt{RETURNS TRIGGER}), sincroniza \texttt{actualizado_en} en cada
    \texttt{UPDATE} sobre \texttt{usuarios} y \texttt{mascotas}, sin que el
    backend deba setear ese campo manualmente.
\end{itemize}

\section{Reproducibilidad y Docker}

ADR-005 formaliza el arranque en un solo comando (\texttt{make up}) y el
pinning de im\'agenes de terceros por digest \texttt{sha256}
(\texttt{postgres:16-alpine}, \texttt{redis:7-alpine} en
\texttt{docker-compose.yml}; \texttt{maven:3.9-eclipse-temurin-21},
\texttt{eclipse-temurin:21-jre-alpine} en \texttt{Backend/Dockerfile}), de
forma que una reconstrucci\'on futura no use, sin darse cuenta, una versi\'on
distinta de una imagen si su \emph{tag} se actualiza silenciosamente en
Docker Hub. El \texttt{Makefile} separa expl\'icitamente operaciones
seguras (\texttt{up}, \texttt{down}, \texttt{clean}, que preservan datos)
de la \'unica operaci\'on destructiva (\texttt{reset-db}, que borra
vol\'umenes, con advertencia expl\'icita antes de ejecutarse).

% EVIDENCIA-FRED-01
% Ruta esperada: figuras/fred/01-docker-healthy.png
\evidencia{figuras/fred/01-docker-healthy.png}%
{Los cuatro servicios (\texttt{postgres}, \texttt{redis}, \texttt{backend}, \texttt{frontend}) en estado \texttt{healthy}/\texttt{started} tras \texttt{make up} desde una clonaci\'on limpia.}%
{fig:fred-docker-healthy}

\section{Redis y cach\'e}

Redis~\cite{redis-docs} respalda tanto la lista negra de JWT como la
cach\'e de lecturas. La cach\'e declarativa de Spring
(\texttt{@Cacheable}/\texttt{@CacheEvict} en \texttt{MascotaService})
respalda las lecturas de \texttt{GET /api/mascotas}. El hit ratio se midi\'o
de dos formas distintas, documentadas ambas por transparencia
metodol\'ogica:

\begin{itemize}[nosep]
  \item \textbf{Global} (\texttt{redis-cli INFO stats},
    \texttt{keyspace_hits}/\texttt{keyspace_misses}, con
    \texttt{CONFIG RESETSTAT} previo): $\approx$49{,}98\,\%. Esta cifra
    cuenta \emph{todas} las lecturas a Redis, incluida la verificaci\'on de
    lista negra de JWT en cada petici\'on autenticada (m\'odulo de Jaime),
    no solo la cach\'e de mascotas --- documentado como hallazgo abierto en
    la primera medici\'on.
  \item \textbf{Aislada} (\texttt{redis-cli DBSIZE} sobre la clave espec\'ifica
    de cach\'e \texttt{mascotas::admin@biopet.ec-0-10-UNSORTED}, estable
    durante una corrida completa de 50~VUs/30~s, sin evicci\'on por memoria
    ni expiraci\'on prematura): cercana al \textbf{100\,\%} tras la primera
    petici\'on de calentamiento, que es la medici\'on m\'as representativa
    del comportamiento real de la cach\'e de mascotas espec\'ificamente.
\end{itemize}

% EVIDENCIA-FRED-05
% Ruta esperada: figuras/fred/05-redis-stats.png
\evidencia{figuras/fred/05-redis-stats.png}%
{\texttt{redis-cli DBSIZE} confirmando una \'unica clave de cach\'e estable durante la corrida de referencia.}%
{fig:fred-redis-stats}

\section{Rendimiento (k6)}
\label{sec:rendimiento-k6}

Herramienta de carga: k6~\cite{k6-docs}.

\begin{table}[H]
  \centering
  \small
  \caption{Resultado real de las seis corridas de k6 contra \texttt{GET /api/mascotas} (HTTPS/TLS~1.3, commit \texttt{a6a8905}). Fuente: \texttt{docs/mediciones/perf/REPORT.md}.}
  \label{tab:k6-resultados}
  \begin{tabular}{@{}lrrrrrrr@{}}
    \toprule
    Corrida & n & Media & Mediana & p95 & p99 & Error & Throughput \\
    & & (ms) & (ms) & (ms) & (ms) & (\%) & (req/s) \\
    \midrule
    run1-caliente & 3233 & 6{,}38 & 5{,}83 & 9{,}41 & 14{,}41 & 0{,}0 & 92{,}14 \\
    run1-fr\'io & 3197 & 12{,}14 & 9{,}39 & 26{,}45 & 44{,}66 & 0{,}0 & 91{,}21 \\
    run2-caliente & 3211 & 9{,}09 & 6{,}39 & 21{,}09 & 59{,}77 & 0{,}0 & 91{,}63 \\
    run2-fr\'io & 3191 & 11{,}89 & 9{,}37 & 27{,}35 & 49{,}01 & 0{,}0 & 91{,}02 \\
    run3-caliente & 3236 & 6{,}45 & 5{,}73 & 10{,}62 & 21{,}08 & 0{,}0 & 92{,}34 \\
    run3-fr\'io & 3226 & 7{,}83 & 6{,}84 & 13{,}74 & 21{,}81 & 0{,}0 & 92{,}05 \\
    \bottomrule
  \end{tabular}
\end{table}

Intervalos de confianza al 95\,\% (distribuci\'on t de Student) para la
media, por ejemplo run1-caliente: $[6{,}18,\ 6{,}58]$~ms; run1-fr\'io:
$[11{,}81,\ 12{,}47]$~ms. La tasa de error fue \textbf{0{,}0\,\%} en las
seis corridas, sin excepci\'on.

% EVIDENCIA-FRED-03
% Ruta esperada: figuras/fred/03-k6-cold.png
\evidencia{figuras/fred/03-k6-cold.png}%
{Salida de consola de una corrida de k6 en fr\'io.}%
{fig:fred-k6-cold}

% EVIDENCIA-FRED-04
% Ruta esperada: figuras/fred/04-k6-warm.png
\evidencia{figuras/fred/04-k6-warm.png}%
{Salida de consola de una corrida de k6 en caliente.}%
{fig:fred-k6-warm}

% EVIDENCIA-FRED-06
% Ruta esperada: figuras/fred/06-performance-report.png
\evidencia{figuras/fred/06-performance-report.png}%
{Vista del reporte agregado de rendimiento (\texttt{docs/mediciones/perf/REPORT.md}).}%
{fig:fred-performance-report}

\section{An\'alisis estad\'istico}

La agregaci\'on de cada corrida (\texttt{scripts/perf-analysis.py}) calcula
media, mediana, desviaci\'on est\'andar muestral, intervalo de confianza al
95\,\% por distribuci\'on t de Student, y los percentiles p50/p90/p95/p99
de \texttt{http_req_duration} directamente sobre los datos crudos de cada
archivo \texttt{k6-runN-\{frio,caliente\}.json}, sin mezclar corridas entre
s\'i.

\section{Limitaciones registradas}

\begin{itemize}[nosep]
  \item Las corridas de k6 se ejecutaron en una \'unica m\'aquina de
    desarrollo, con duraci\'on limitada ($\sim$30--35~s por corrida, 50~VUs);
    no se evalu\'o comportamiento bajo cargas sostenidas de mayor duraci\'on
    ni con m\'as usuarios virtuales.
  \item El hit ratio ``global'' de Redis ($\approx$49{,}98\,\%) mezcla
    lecturas de cach\'e con lecturas de lista negra de JWT; se documenta la
    medici\'on aislada como la m\'as representativa, sin ocultar la
    primera cifra.
  \item La contrase\~na de desarrollo de \texttt{biopet_app} queda en texto
    plano en \texttt{db/roles.sql} y \texttt{.env.example} porque los
    scripts de inicializaci\'on de Postgres no pueden leer variables de
    entorno de Spring Boot; se documenta como valor exclusivamente de
    desarrollo (ADR-004).
  \item \texttt{db/schema.sql} es una copia manual de
    \texttt{V1__schema_inicial.sql}: si se agrega una nueva migraci\'on
    Flyway, debe actualizarse manualmente para no quedar desfasada en una
    clonaci\'on limpia.
  \item El trigger \texttt{set_actualizado_en} no cuenta con un test de
    integraci\'on dedicado exclusivamente a \'el; se verifica de forma
    impl\'icita en las pruebas que hacen \texttt{UPDATE}.
\end{itemize}
